\chapter{Creutzfeldt-Jakob Disease (CJD)}
\index{Creutzfeldt-Jakob Disease (CJD)}
\label{sec:Creutzfeldt-Jakob Disease}

\section{Overview}
Creutzfeldt-Jakob disease is a rapidly progressive, fatal neurodegenerative disease caused by prions---misfolded, protease-resistant isoforms (PrP$^{Sc}$) of the normal cellular prion protein (PrP$^{C}$), with an incidence of 1--2 per million annually worldwide, median age at onset 60--65 years, with sporadic CJD (sCJD) accounting for 85--90\% of cases, familial CJD (10\%) linked to PRNP mutations, and iatrogenic CJD ($<$1\%) from contaminated surgical instruments.
\section{Causes}
This condition happens because of PrP$^{Sc}$ converting native PrP$^{C}$ into the misfolded conformation through a template-directed process, leading to exponential accumulation of PrP$^{Sc}$ aggregates, causing vacuolation (spongiform change), neuronal loss, and astrocytosis. The underlying mechanisms involve complex interactions between genetic predisposition and environmental factors that disrupt normal physiological processes. The underlying mechanisms involve complex interactions between genetic predisposition and environmental factors. Disruption of normal physiological processes leads to the characteristic features of the condition. Multiple contributing factors may act synergistically to trigger or worsen the condition. Understanding the specific cause helps guide targeted treatment approaches.
\section{Risk Factors}


\begin{itemize}[noitemsep]
  \item Genetic predisposition and family history
  \item Advancing age
  \item Lifestyle factors (diet, exercise, smoking, alcohol)
  \item Environmental exposures
  \item Underlying medical conditions
  \item Medication use
\end{itemize}
\section{Signs and Symptoms}

\subsection{Common symptoms}
\begin{itemize}[noitemsep]
  \item Symptoms vary depending on the severity and extent of the condition Pain or discomfort in the affected area Functional impairment affecting daily activities Changes in appearance or function of affected tissues Systemic symptoms may include fatigue, fever, or weight changes Symptoms may develop gradually or appear suddenly
  \item Pain or discomfort in the affected area
  \end{itemize}

\subsection{Emergency warning signs}
\begin{itemize}[noitemsep]
  \item Severe or worsening symptoms of creutzfeldt-jakob disease (cjd) require immediate medical evaluation
\end{itemize}
\section{Progression and Stages}
You may experience rapidly progressive dementia (weeks to months), myoclonus (startle-sensitive, characteristic), cerebellar loss of coordination, pyramidal and extrapyramidal signs, visual disturbances, and akinetic mutism in the terminal stage, with death typically occurring within 1 year of symptom onset. The condition may progress through different stages of severity over time. Early stages often involve mild or subtle symptoms that may be overlooked. Moderate stages involve more noticeable symptoms affecting daily function. Advanced stages may involve significant impairment and complications.
\section{Complications}
\begin{itemize}[noitemsep]
  \item Disease progression leading to worsening symptoms
  \item Secondary infections or injuries
  \item Side effects of treatment
  \item Reduced quality of life
  \item Emotional and psychological impact
\end{itemize}
\section{Diagnosis}
Doctors diagnose this using clinical criteria, EEG showing periodic sharp wave complexes (PSWCs), CSF RT-QuIC (specific and sensitive), and MRI showing T2/FLAIR hyperintensity in the basal ganglia and thalamus. Comprehensive medical history and physical examination Laboratory tests to assess organ function and rule out other conditions Imaging studies to visualize affected structures Specialized testing depending on the affected system Consultation with relevant specialists Monitoring of disease progression over time

\begin{itemize}[noitemsep]
  \item Comprehensive medical history and physical examination
  \item Laboratory tests to assess organ function
  \item Imaging studies to visualize affected structures
  \item Specialized testing depending on the affected system
  \item Consultation with relevant specialists
\end{itemize}
\section{Prevention}
\begin{itemize}[noitemsep]
  \item Maintain a healthy lifestyle with balanced diet and exercise
  \item Avoid known risk factors
  \item Attend regular medical check-ups for early detection
  \item Follow recommended screening guidelines
\end{itemize}
\section{Treatment Options}
Medications to manage symptoms and address underlying mechanisms Lifestyle modifications and self-care measures Physical therapy and rehabilitation Surgical intervention when indicated Regular monitoring and follow-up care Patient education and support services

\begin{itemize}[noitemsep]
  \item Medications to manage symptoms and address underlying mechanisms
  \item Lifestyle modifications and self-care measures
  \item Physical therapy and rehabilitation
  \item Surgical intervention when indicated
  \item Regular monitoring and follow-up care
\end{itemize}
\section{Living with It}
Treatment typically involves palliative.

\begin{itemize}[noitemsep]
  \item Follow your treatment plan as prescribed
  \item Maintain regular follow-up with your healthcare provider
  \item Make necessary lifestyle adjustments
  \item Seek support from family, friends, or support groups
  \item Stay informed about your condition
\end{itemize}
\section{Outlook}
outlook is uniformly fatal. The outlook varies depending on the specific condition, severity at diagnosis, and response to treatment. Early intervention generally leads to better outcomes. Many conditions can be effectively managed with appropriate treatment. Regular follow-up care is important for monitoring and treatment adjustment.
